\setlength{\footskip}{8mm}

\chapter{Literature Review}
\label{ch:literature-review}

\section{Overview}
\label{sec:lit-overview}

This chapter surveys the literature underlying the architecture
developed in this thesis, organized to trace the field's evolution from
its earliest formulations to the specific research gap this work
addresses. Section~\ref{sec:lit-pipeline} establishes the two-stage
2D-to-3D lifting pipeline that frames the rest of the chapter.
Section~\ref{sec:lit-conv-graph} reviews convolutional and graph-based
lifting approaches, including early treatments of the human skeleton as a
graph. Section~\ref{sec:lit-transformer} traces the Transformer-based
lineage of 3D human pose estimation methods and the quadratic-complexity
limitation they share. Section~\ref{sec:lit-ssm-foundations} introduces
structured state space models and the Mamba architecture as a
linear-complexity alternative to self-attention.
Section~\ref{sec:lit-ssm-pose} reviews the emerging body of work applying
state space models to human motion and pose estimation directly.
Section~\ref{sec:lit-skeleton-ordering} examines how sequential models
represent and order graph-structured data such as the human skeleton,
drawing on evidence from outside the pose estimation literature. Finally,
Section~\ref{sec:lit-summary-gap} synthesizes this review into an explicit
statement of the research gap, connecting directly to the problem
statement in Section~\ref{sec:intro-problem-statement} and motivating the
architecture presented in Chapter~\ref{ch:methodology}.

\section{The Two-Stage 3D Human Pose Estimation Pipeline}
\label{sec:lit-pipeline}

The majority of monocular 3D human pose estimation methods, including the
architecture developed in this thesis, decompose the task into two
sequential stages: 2D pose detection, followed by 2D-to-3D lifting.

In the first stage, a 2D pose detector processes an image or video frame
and localizes a fixed set of anatomical keypoints (e.g., wrists, elbows,
knees) on the image plane. HRNet \cite{sun2019hrnet} is a representative
and widely adopted architecture for this stage: unlike earlier
encoder-decoder designs that recover resolution only at the output,
HRNet maintains high-resolution feature representations in parallel with
progressively lower-resolution ones throughout the network, repeatedly
exchanging information across resolutions. This design yields more
spatially precise keypoint localization, which in turn provides a more
reliable input signal to the second stage. As established in
Section~\ref{sec:intro-limitations}, this thesis treats 2D keypoint
detection as an external, fixed input rather than a component to be
modified; HRNet-detected keypoints (or ground-truth 2D annotations, where
available) are used directly as the input representation.

The second stage, 2D-to-3D lifting, takes the resulting 2D keypoint
sequence and infers the corresponding 3D joint coordinates. Formally,
given a sequence of $T$ frames of 2D keypoints $X \in
\mathbb{R}^{T \times J \times 2}$ for $J$ body joints, a lifting network
learns a mapping to 3D joint coordinates $Y \in \mathbb{R}^{T \times J
\times 3}$ (or, in single-frame-output formulations, to the 3D pose of a
single target frame within the input window). VideoPose3D
\cite{pavllo2019videopose3d} formalized this sequence-to-pose (or
sequence-to-sequence) lifting formulation using dilated temporal
convolutions over the 2D keypoint sequence, establishing both the general
problem setup and the practice of exploiting multiple input frames to
resolve the depth ambiguity and self-occlusion inherent to any single 2D
projection. This two-stage, sequence-based lifting formulation, together
with its associated evaluation protocol on the Human3.6M
\cite{ionescu2014human36m} and MPI-INF-3DHP \cite{mehta2017mpiinf3dhp}
benchmarks using Mean Per Joint Position Error (MPJPE), remains the
standard framework within which nearly all subsequent lifting
architectures, including the one proposed in this thesis, are designed and
evaluated.

What distinguishes the methods reviewed in the remainder of this chapter
is not this overall pipeline, which they share, but the architecture used
to model spatial (inter-joint) and temporal (inter-frame) dependencies
within the lifting network itself. Section~\ref{sec:lit-conv-graph}
begins this comparison with convolutional and graph-based designs.

\section{Convolutional and Graph-Based Approaches to Pose Lifting}
\label{sec:lit-conv-graph}

VideoPose3D \cite{pavllo2019videopose3d}, introduced in
Section~\ref{sec:lit-pipeline}, models temporal dependencies across frames
through dilated 1D convolutions, but treats the $J$ joints within a frame
as independent channels: it does not explicitly encode which joints are
anatomically connected to which. A parallel line of work addressed this by
representing the human body directly as a graph, in which nodes correspond
to joints and edges correspond to bones, and applying graph convolution to
propagate information along this structure.

Spatial Temporal Graph Convolutional Networks (ST-GCN) \cite{yan2018stgcn}
established this idea for skeleton-based recognition, defining spatial
graph convolutions over the fixed, physical joint-adjacency structure of
the human body and combining them with temporal convolutions across
frames. Although originally developed for action recognition rather than
3D pose regression, ST-GCN's core contribution, treating the skeleton
as a graph with an explicit, anatomically-defined adjacency rather than
as an unordered or arbitrarily-ordered set of joints, directly
underlies the graph- and topology-aware pose estimation methods that
followed, and is foundational to the kinematic-tree perspective adopted
later in this thesis (Section~\ref{sec:lit-skeleton-ordering}).

Subsequent work observed that a fixed physical adjacency is restrictive:
it only allows information to propagate between directly-connected joints
in a single graph-convolution layer, and cannot represent relationships
that are semantically meaningful but not anatomically adjacent (e.g.,
between the two hands during a clapping motion). Semantic Graph
Convolutional Networks (SemGCN) \cite{zhao2019semgcn} addressed this by
learning edge weights end-to-end from data rather than fixing them to the
physical skeleton, allowing the network to discover task-relevant joint
relationships beyond the kinematic chain. GLA-GCN \cite{yu2023glagcn}
extended this further for video-based 3D HPE specifically, combining a
global representation of the whole skeleton with local representations of
joint groups (e.g., individual limbs), and aggregating both to produce the
final pose estimate. This global-local decomposition reappears, in a
different architectural form, in the state-space-based methods reviewed in
Section~\ref{sec:lit-ssm-pose}.

Despite these refinements, graph convolutional approaches share a common
limitation: each graph-convolution layer only aggregates information from
a node's immediate neighborhood, so capturing dependencies between distant
joints or distant frames requires stacking many layers, which is
parameter-inefficient and can dilute long-range signal. This limitation,
shared with the temporal convolutions of VideoPose3D, motivated a shift
toward architectures with an inherently global receptive field, namely
the self-attention mechanism reviewed next.

\section{Transformer-Based 3D Human Pose Estimation}
\label{sec:lit-transformer}

The Transformer architecture \cite{vaswani2017attention} replaces
recurrence and convolution with self-attention, in which every element of
a sequence directly attends to every other element in a single layer.
This gives self-attention an inherently global receptive field, in
contrast to the local receptive fields of the convolutional and
graph-based methods reviewed in Section~\ref{sec:lit-conv-graph}. Following
its success in language modeling and, via the Vision Transformer
\cite{dosovitskiy2020vit}, in image classification, self-attention was
adopted for 3D human pose estimation to model spatial (inter-joint) and
temporal (inter-frame) dependencies without the depth or layer-count
restrictions of convolution and graph convolution.

PoseFormer \cite{zheng2021poseformer} was the first method to apply this
idea to 3D HPE, using a spatial self-attention block to model relationships
among joints within a frame and a temporal self-attention block to model
relationships across frames, estimating the 3D pose of a center frame from
a surrounding temporal window. Subsequent work extended this design along
several complementary directions. MHFormer \cite{li2022mhformer} generates
multiple plausible 3D pose hypotheses in parallel and aggregates them
end-to-end, directly targeting the depth ambiguity inherent to monocular
input. MixSTE \cite{zhang2022mixste} alternates spatial and temporal
transformer blocks within a sequence-to-sequence formulation, estimating
3D poses for every frame in the input window rather than only the center
frame, which improves temporal coherence across the output sequence.
PoseFormerV2 \cite{zhao2023poseformerv2} instead targets efficiency and
robustness to noisy 2D input by transforming part of the temporal sequence
into the frequency domain, reducing the effective sequence length
processed by self-attention without a proportional loss of temporal
information.

STCFormer \cite{tang2023stcformer} argued that modeling spatial and
temporal correlations as two separate, sequential attention operations
(as in PoseFormer and its successors) is suboptimal, and instead models
them jointly through a criss-cross attention mechanism operating over both
axes simultaneously. MotionBERT \cite{zhu2023motionbert} shifted the focus
from architecture to representation, pretraining a unified motion encoder
on large-scale 2D-to-3D data that can subsequently be fine-tuned for pose
estimation and other motion-related tasks. MotionAGFormer
\cite{mehraban2024motionagformer} combined a Transformer branch with a
graph-convolutional branch (a "GCNFormer") in parallel, explicitly
recombining the global modeling capacity of self-attention with the
explicit skeletal-structure bias of graph convolution reviewed in
Section~\ref{sec:lit-conv-graph}, and represents one of the strongest
Transformer-based baselines against which the architecture proposed in
this thesis is compared in Chapter~\ref{ch:results}.

In parallel, efficiency-oriented Transformer variants such as the Swin
Transformer \cite{liu2021swin} attempted to reduce self-attention's
computational cost by restricting attention to local, shifted windows
rather than the full sequence, trading some global context for
near-linear scaling. While effective in general vision tasks, this
windowing strategy has not displaced full self-attention as the dominant
choice in 3D HPE, where global spatio-temporal context is directly tied to
resolving depth ambiguity.

Across this entire lineage, the underlying limitation identified in
Section~\ref{sec:intro-problem-statement} persists: self-attention's
computational and memory cost scales quadratically with sequence length,
regardless of whether that sequence is processed jointly, alternately, in
the frequency domain, or in restricted windows. This shared constraint on
scaling temporal context motivated the search for an architecture that
retains a global receptive field while scaling linearly with sequence
length: the state space model, reviewed next.

\section{State Space Models: From S4 to Mamba}
\label{sec:lit-ssm-foundations}

Structured state space models (SSMs) trace back to HiPPO
\cite{gu2020hippo}, which framed the problem of a network remembering its
input history as one of online function approximation, and derived a
memory-update mechanism that optimally compresses a continuous or discrete
sequence's history onto a fixed-size state by projection onto a
polynomial basis. Building on this theoretical foundation, the Linear
State-Space Layer (LSSL) \cite{gu2021lssl} showed that a single linear
state-space formulation, $\dot{x} = Ax + Bu,\; y = Cx + Du$, can be viewed
simultaneously as a recurrent model, a convolutional model, and a
continuous-time (neural ODE) model, unifying three previously separate
families of sequence models and inheriting their respective strengths:
constant-memory recurrent inference, parallelizable convolutional
training, and principled continuous-time behavior.

The structured state space sequence model (S4) \cite{gu2021s4} made this
formulation practical at scale by imposing a structured (diagonal-plus-low-rank)
parameterization on the state matrix $A$, sharply reducing the
computational cost of the convolutional view while preserving HiPPO's
long-range memory properties, and demonstrated strong performance on
tasks requiring modeling of very long sequences that had previously been
difficult for both RNNs and Transformers. S5 \cite{smith2022s5} further
simplified this design by replacing S4's bank of independent single-input,
single-output SSMs with a single multi-input, multi-output SSM computed
via a parallel scan, reducing implementation complexity while retaining
comparable performance.

Mamba \cite{gu2023mamba} introduced a further refinement: rather than
fixing the SSM's parameters for the entire sequence, Mamba makes them
\emph{selective}, computing them as functions of the current input, which
lets the model dynamically decide what information to retain or discard as
it processes a sequence, a capability that fixed-parameter SSMs such as S4
lack. Combined with a hardware-aware parallel scan
algorithm, Mamba achieves Transformer-competitive quality across language
and other sequence modeling benchmarks while retaining linear-time and
linear-memory complexity in sequence length, directly addressing the
quadratic-complexity limitation identified in
Section~\ref{sec:lit-transformer}.

Because Mamba's selective scan was originally formulated for causal,
unidirectional sequences (as is natural for autoregressive language
modeling), applying it to non-causal vision data required further
adaptation. Vision Mamba \cite{zhu2024visionmamba} addressed this by
processing an image's flattened patch sequence in both forward and
backward directions and combining the two, allowing each patch to be
informed by patches on both sides rather than only those preceding it in
the flattened order. VMamba \cite{liu2024vmamba} extended this further
with a cross-scan module that traverses a 2D feature map along multiple
directions (e.g., row-major and column-major, in both directions) to
better preserve spatial adjacency that a single, fixed 1D flattening would
otherwise disrupt. These two designs sit within a fast-growing body of
vision and general state space models, now large enough to warrant dedicated
surveys \cite{zhang2024visualmambasurvey, qu2024mambasurvey}. Both works
establish a principle central to this thesis: when a state space model's strictly sequential inductive bias is
applied to data that is not inherently sequential, \emph{how} that data is
linearized into a sequence, the scan order, is not incidental but a
first-class architectural design choice with direct consequences for what
the model can represent. Section~\ref{sec:lit-ssm-pose} examines how this
principle has, so far, been addressed for human motion and pose data
specifically.

\section{State Space Models for Human Motion and Pose Estimation}
\label{sec:lit-ssm-pose}

Following the adoption of state space models in general vision tasks
(Section~\ref{sec:lit-ssm-foundations}), recent work has applied SSMs to
human motion and pose estimation directly. The reach of this idea extends
beyond pose regression: Motion Mamba \cite{zhang2024motionmamba} builds an
SSM backbone for long-sequence human motion generation, and Simba
\cite{chaudhuri2024simba} pairs a Mamba module with a shift graph
convolution for skeleton-based action recognition, both showing that an SSM
can model the dynamics of articulated human movement at a length that strains
attention. Within pose estimation, HumMUSS \cite{mondal2024hummuss} was among
the first to apply state space modeling to human motion understanding, and
Pose Magic \cite{zhang2024posemagic} combined a Mamba backbone with a graph
convolutional component to keep an explicit skeletal-structure bias alongside
the linear-complexity temporal modeling of Mamba.

The most direct precedent for the present work is PoseMamba
\cite{huang2025posemamba}, a pure-Mamba network for 2D-to-3D lifting built
from bidirectional global-local spatio-temporal blocks. PoseMamba scans the
joints of a frame in their index order for a global view and adds a second,
locally-grouped scan to recover fine skeletal detail, a reordering it
motivates as giving the state space model a more sensible geometric order
than a single index sweep. This is the same observation the present work
starts from, that the spatial scan order matters, but PoseMamba's reordering
is a hand-designed grouping rather than an order read from the skeleton's
kinematic tree.

Across this body of work, echoing the discussion in
Section~\ref{sec:lit-ssm-foundations}, a plain unidirectional SSM scan
applied naively to a flattened skeleton sequence underperforms, and some
explicit treatment of the skeleton's non-sequential structure proves
necessary. Different methods address this in different ways. Skeleton
Mamba \cite{nguyen2025skeletonmamba} introduces Group Scan and Joint Scan
strategies that learn how to group and order joints for the SSM's spatial
scan, alongside a separate temporal scan across frames. Mamba-Driven
Topology Fusion \cite{zheng2025mambatopology} instead fuses explicit
skeleton-topology priors with the SSM's sequence representation, rather
than relying on the scan order alone to convey structural information. A
Structure-Aware and Motion-Adaptive Framework \cite{lu2025structureawaremamba}
combines a structure-aware skeleton encoding with motion-adaptive SSM
blocks that adjust their behavior based on the motion dynamics of the
input sequence. High-Resolution Spatiotemporal Modeling with Global-Local
State Space Models \cite{feng2025highresssm} extends the global-local
spatial decomposition seen earlier in GLA-GCN
(Section~\ref{sec:lit-conv-graph}) to high-resolution, video-based pose
estimation within an SSM framework. More recently, SasMamba
\cite{cui2025sasmamba} proposes a lightweight, structure-aware stride SSM
that reduces effective sequence length while preserving accuracy, and
DBMambaPose \cite{wang2025dbmambapose} decouples the spatial and temporal
bidirectional SSM computations, arguing that processing them jointly, as
in earlier global-local designs, is not necessary for strong performance.
Closest in spirit to an anatomy-aware design is HGMamba
\cite{cui2025hgmamba}, which runs a hypergraph convolution branch for local
joint structure in parallel with a Mamba branch for global spatio-temporal
context and fuses the two. HGMamba adds structure through a separate
graph-convolution branch; the present work instead injects the skeletal
structure into the Mamba scan itself, through the order in which joints are
visited, leaving the backbone otherwise unchanged.

Across this body of work, three broad strategies for handling the
skeleton's non-sequential structure within an SSM recur: (i) learning the
scan order or grouping directly from data (e.g., Group/Joint Scan), (ii)
fusing separately-computed topology or structure information alongside the
sequence representation, and (iii) decoupling or restructuring the
bidirectional scan computation itself. What none of these strategies
provides is a scan order derived \emph{directly}
from the skeleton's own kinematic hierarchy, one that uses the
parent-child joint structure of the body itself, rather than a learned
grouping or an auxiliary topology-fusion module, to determine the order in
which joints are presented to the SSM. Section~\ref{sec:lit-skeleton-ordering} examines
the representational basis for such a kinematic-tree-derived ordering.

\section{Skeleton Representation and Sequence Ordering for Sequential Models}
\label{sec:lit-skeleton-ordering}

This section draws together three lines of evidence, largely from outside
the pose estimation literature itself, that together motivate a
kinematic-tree-derived scan order for sequential models such as the
state space models reviewed in Section~\ref{sec:lit-ssm-foundations}.

\textbf{The human body has a formal, well-defined kinematic tree
structure.} As established in the parametric body modeling literature,
SMPL \cite{loper2015smpl} represents the human body's pose using a
standard skeletal rig in which each joint's rotation is defined relative
to its parent joint in a kinematic tree rooted at the pelvis. This
parent-child hierarchy is not specific to SMPL: the same structure
underlies the joint definitions of the benchmark datasets used
throughout this thesis (Section~\ref{sec:lit-pipeline}) and the
skeleton-as-graph representations reviewed in
Section~\ref{sec:lit-conv-graph} (e.g., ST-GCN's fixed physical adjacency
\cite{yan2018stgcn}). A rooted tree, unlike a general graph, admits a
canonical traversal order in a way that a general graph does not: starting
from the root, nodes can be visited layer by layer, with every node
visited only after its parent.

\textbf{Sequential models benefit from context in both directions when
the underlying data is not strictly causal.} Long Short-Term Memory (LSTM)
networks \cite{hochreiter1997lstm} established gated recurrent memory as
the standard mechanism for sequential modeling, but process input in a
single direction. BERT \cite{devlin2019bert} demonstrated that, for data
without an inherent causal or temporal ordering, processing a sequence
bidirectionally and allowing each element to attend to context on both
sides improves the resulting representation. This principle
directly underlies the bidirectional scanning strategies of Vision Mamba
and VMamba (Section~\ref{sec:lit-ssm-foundations}), and, by extension, of
the bidirectional spatio-temporal architecture developed in this thesis: a
skeleton's spatial joint arrangement, like an image's spatial layout, has
no inherent causal direction, and so a joint's representation should be
informed by structural context in both directions along the traversal, not
only by joints visited before it.

\textbf{A principled node ordering measurably improves
sequential modeling of graph-structured data.} Outside the pose estimation
and vision literature, GraphRNN \cite{you2018graphrnn} studied the problem
of generating graphs using autoregressive sequence models, and found that
restricting the node ordering used to linearize a graph into a sequence to
a breadth-first search (BFS) ordering markedly improved both the
scalability and the quality of the resulting sequence model, compared to
an arbitrary or random node ordering. The intuition underlying this result
generalizes directly to the problem addressed in this thesis: a BFS
ordering keeps nodes that are close in the graph close in the resulting
sequence, which is precisely the locality property that benefits a
sequence model's ability to capture local dependencies, as discussed for
SSMs in Section~\ref{sec:lit-ssm-foundations}. Later work reinforces the
point. A dedicated study of ordering in autoregressive graph generation
\cite{bu2023ordering} compares breadth-first, depth-first, and canonical
orderings and confirms that the choice of node order has a measurable effect
on what the sequence model learns. The question has resurfaced with state
space models specifically: applying a Mamba scan to a graph requires
serializing its nodes, and both Graph Mamba \cite{behrouz2024graphmamba} and
Graph-Mamba \cite{wang2024graphmamba} identify the absence of a canonical
node order as the central obstacle, addressing it with node prioritization
and permutation schemes. A human skeleton avoids that obstacle, because a
rooted kinematic tree already supplies a canonical breadth-first order.

These three observations together suggest a specific, previously
unexplored design choice for SSM-based pose estimation, distinct from the
learned-scan and topology-fusion strategies reviewed in
Section~\ref{sec:lit-ssm-pose}: since the human skeleton is formally a
rooted kinematic tree (SMPL), since bidirectional context benefits
non-causal structural data (BERT, Vision Mamba, VMamba), and since BFS
ordering is a principled way to linearize a graph for a sequence model
while preserving locality (GraphRNN), a breadth-first
traversal of the skeleton's kinematic tree, combined with bidirectional
scanning, offers an anatomically-grounded alternative to both fixed
index-based orderings and learned scan strategies. Section~\ref{sec:lit-summary-gap}
formalizes this observation as the specific research gap this thesis
addresses.

\section{Summary and Research Gap}
\label{sec:lit-summary-gap}

This chapter traced the evolution of monocular 3D human pose estimation
across three architectural paradigms. Convolutional and graph-based
lifting networks (Section~\ref{sec:lit-conv-graph}) introduced explicit
skeletal structure but are limited by local receptive fields.
Transformer-based methods (Section~\ref{sec:lit-transformer}) replaced
this local receptive field with global self-attention, achieving strong
accuracy at the cost of quadratic computational complexity in sequence
length. State space models (Sections~\ref{sec:lit-ssm-foundations}
and~\ref{sec:lit-ssm-pose}) offer a resolution to this cost, matching
Transformer-level modeling capacity with linear complexity, and have
recently been applied to human motion and pose estimation directly.

Across this final paradigm, however, a consistent gap remains. Every
SSM-based pose estimation method reviewed in
Section~\ref{sec:lit-ssm-pose} must confront the same underlying
mismatch identified in Section~\ref{sec:intro-problem-statement}: an SSM
processes a strictly one-dimensional sequence, while the human skeleton is
a hierarchical, tree-structured graph with no canonical linear order. The
solutions proposed so far (learned scan groupings, auxiliary topology
fusion, or restructured bidirectional computation) address this mismatch
without deriving the scan order directly from the skeleton's own
kinematic structure. Yet Section~\ref{sec:lit-skeleton-ordering} showed
that the necessary ingredients for such a derivation already exist
separately in the literature: a formal kinematic-tree definition of the
body (SMPL), evidence that bidirectional context benefits non-causal
structural data (BERT, Vision Mamba, VMamba), and evidence that a
breadth-first traversal is a principled, locality-preserving way to
linearize a graph for a sequence model (GraphRNN). No existing SSM-based
pose estimation method, however, combines them into a
kinematic-tree-derived scan order.

This is the specific research gap this thesis addresses. Chapter~\ref{ch:methodology}
presents an architecture aimed at this gap: a bidirectional
spatio-temporal state space backbone, consistent with the efficiency
motivation established in this chapter, in which the spatial scan order is
given by a breadth-first search traversal of the skeleton's kinematic
tree rather than a fixed index order or a learned grouping.

\FloatBarrier

